% !TEX root = main.tex
% !TEX encoding = UTF-8 Unicode

\section{Demonstration Objective and Success Criteria}

The demonstration should prove three claims:
\begin{itemize}
  \item The chosen method (SCRFD) is technically justified by both chapter and paper.
  \item The system can run in practical real-time settings.
  \item The accuracy-efficiency trade-off is visible, measurable, and reproducible.
\end{itemize}

\section{What We Will Deliver as the SOTA Demonstration}

\subsection{Model set and expected role}
\begin{itemize}
  \item SCRFD-0.5GF: ultra-light deployment profile.
  \item SCRFD-2.5GF: balanced profile for the main demonstration.
  \item SCRFD-10GF or 34GF (if hardware permits): high-accuracy profile.
\end{itemize}

\subsection{Evaluation conditions}
\begin{itemize}
  \item Input mode A: VGA single-scale stream for speed-oriented comparison.
  \item Input mode B: original-resolution images for quality stress testing.
  \item Input mode C: tiny-face crowded scenes to test hard-case robustness.
\end{itemize}

\section{Demonstration Protocol (Step-by-Step)}

\begin{enumerate}
  \item Show baseline setup (hardware, input resolution, model profile).
  \item Run SCRFD in real-time and display FPS, latency, and detections.
  \item Compare at least two model sizes to show trade-off trend.
  \item Show tiny-face or cluttered scenes and discuss observed failure modes.
  \item Summarize results in one compact comparison table.
\end{enumerate}

\section{Metrics and Artifacts to Prepare for TAs}

\subsection{Mandatory metrics}
\begin{itemize}
  \item Latency per frame (ms).
  \item Throughput (FPS).
  \item Detection quality indicator (AP if available, otherwise labeled subset precision/recall).
  \item Parameter count and GFLOPs per model profile.
\end{itemize}

\subsection{Artifacts to submit/show}
\begin{itemize}
  \item One reproducible command list for each demo mode.
  \item One summary table with all compared settings.
  \item One short demo video (fallback if live demo is unstable).
  \item One page of qualitative examples (easy, medium, hard cases).
\end{itemize}

\section{How to Give TAs Enough Technical Details}

\begin{itemize}
  \item Clearly separate \textit{paper-reported numbers} and \textit{team-reproduced numbers}.
  \item Provide exact model variant names, input size, and runtime environment.
  \item Explain every major design choice: why this model size, why this resolution, why this metric.
  \item Report at least one failure case and proposed fix, not only successful examples.
\end{itemize}

\section{Oral Presentation Plan}

\begin{table}[ht]
\centering
\caption{Suggested 10-minute oral structure}
\label{tab:oral-plan}
\begingroup
\renewcommand{\arraystretch}{1.18}
\small
\setlength{\tabcolsep}{4pt}
\begin{tabularx}{\textwidth}{>{\RaggedRight\arraybackslash}p{0.12\linewidth} >{\RaggedRight\arraybackslash}X >{\RaggedRight\arraybackslash}p{0.20\linewidth}}
\toprule
\textbf{Time} & \textbf{Content} & \textbf{Expected output} \\
\midrule
0:00--1:30 & Problem, chapter scope, and paper scope & Clear framing and motivation \\
1:30--4:00 & Textbook method and SCRFD method summary & Technical understanding \\
4:00--6:30 & Chapter-paper relationship and critical comparison & Methodological depth \\
6:30--8:30 & Live/recorded demo and quantitative table & Practical validation \\
8:30--10:00 & Limitations, extension plan, and Q\&A setup & Research maturity and extensibility \\
\bottomrule
\end{tabularx}
\endgroup
\end{table}

\section{Execution Timeline to Reach Submission}

\begin{enumerate}
  \item Finalize report text and figures.
  \item Prepare reproducible demo scripts and metric extraction.
  \item Rehearse oral flow with strict timing.
  \item Freeze final PDF and backup demo video.
\end{enumerate}
